\chapter{Wymagania}\label{ch:wymagania}

% TODO: formalizacja wymagań — ok. 5--8 stron.

\section{Interesariusze}\label{sec:interesariusze}

% TODO: lekarze POZ, specjaliści, pacjenci, administrator placówki, inspektor
% ochrony danych osobowych.

\section{Wymagania funkcjonalne}\label{sec:wymagania-funkcjonalne}

% TODO: lista numerowana F1, F2, ... z krótkim opisem każdego wymagania.
\begin{enumerate}
  \item[\textbf{F1}] System rejestruje przebieg wizyty lekarskiej w~formie
        nagrania audio.
  \item[\textbf{F2}] System transkrybuje nagranie do~tekstu z~rozróżnieniem
        mówców (lekarz / pacjent).
  \item[\textbf{F3}] System generuje notatkę SOAP (Subjective, Objective,
        Assessment, Plan) na~podstawie transkrypcji.
  \item[\textbf{F4}] System anonimizuje wszelkie dane osobowe pacjenta przed
        wysłaniem danych do~chmury.
  \item[\textbf{F5}] System sugeruje leki refundowane na~podstawie bazy NFZ.
  \item[\textbf{F6}] Lekarz akceptuje lub odrzuca każdą sugestię przed
        zapisaniem w~dokumentacji.
\end{enumerate}

\section{Wymagania niefunkcjonalne}\label{sec:wymagania-niefunkcjonalne}

% TODO: wydajność, bezpieczeństwo, zgodność z~RODO, dostępność, skalowalność.
\begin{enumerate}
  \item[\textbf{N1}] \textbf{Bezpieczeństwo danych:} surowe audio i~niezanonimizowana
        transkrypcja nie~mogą opuszczać infrastruktury placówki (strefa
        zaufania lokalna).
  \item[\textbf{N2}] \textbf{Zgodność z~RODO:} każda operacja na~danych
        osobowych jest logowana (audytowalność).
  \item[\textbf{N3}] \textbf{Czas odpowiedzi:} generacja notatki SOAP dla
        15-minutowej wizyty w~czasie $\leq 3$~minut.
  \item[\textbf{N4}] \textbf{Dokładność transkrypcji:} WER $\leq 15\%$ dla
        próbki testowej dialogów medycznych w~języku polskim.
\end{enumerate}

\section{Przypadki użycia}\label{sec:use-cases}

% TODO: diagramy UML przypadków użycia, scenariusze główny + alternatywne.
